\section{The linear-field bound and simplex rigidity}\label{sec:06-linear-equality}

Equation~\eqref{eq:23} leaves us to estimate the terms depending on the
linear Gaussian field $L=\sum_{\ell=1}^m G_\ell M_\ell$ and the fixed
matrix $T$. We now evaluate these terms and use the segment
inequality~\eqref{eq:4} to control their sum.  The strict terms retained in this calculation will also determine
the equality case.

\subsection{Spectral averaging and Gaussian integration by parts}

The matrices \(T\) and \(M_\ell\) are fixed throughout this section, as in
\eqref{eq:8}.  At a given Gaussian input, choose an orthonormal eigenbasis of
\(L\), with eigenvalues \(x_1,\ldots,x_m\), and express all matrix entries
in that basis.  The indices \(i,j\) refer to this spectral basis;
\(\ell\) still labels the fixed coefficient \(M_\ell\) and the Gaussian
variable \(G_\ell\). Only the matrix coordinates are changed.
Define a positive symmetric measure \(\nu\) on
\(\mathbb R^2\) by
\[
 \nu[\Phi]=\frac1m\mathbb E\sum_{\ell,i,j}
       |(M_\ell)_{ij}|^2\Phi(x_i,x_j),
 \qquad \mathbf S=\sum_\ell M_\ell^2.
\]
At each Gaussian input, it assigns weight
\(m^{-1}\sum_\ell|(M_\ell)_{ij}|^2\) to the ordered pair
\((x_i,x_j)\), and then averages over that input. It is not in general
a probability measure.
This definition is independent of the chosen eigenbasis: between two
fixed eigenspaces, the sum of the squared entries is the squared
Hilbert--Schmidt norm of the corresponding block.  In particular, repeated
eigenvalues cause no ambiguity.  The measure has total mass
\(\nu[1]=\tau\mathbf S<\infty\), and it integrates every function of
polynomial growth, because \(L\) is linear in a finite-dimensional Gaussian.

For a scalar or vector profile \(f\), write
\[
 \bar f_{ij}=\int_0^1 f((1-t)x_i+tx_j)\,dt.
\]
Thus \(\bar f_{ij}=f(x_i)\) when \(x_i=x_j\).  For a scalar profile define
\[
 \mathfrak R(f)=\sum_\ell\mathbb E\left[
       M_\ell\circ(\bar f\odot M_\ell)-M_\ell^2\circ f(L)\right].
\]
Here \(\odot\) denotes entrywise multiplication in an eigenbasis of \(L\).
The resulting matrix is returned to the fixed coordinates before taking
expectation.  This convention is essential: the matrices \(M_\ell\) are
deterministic in the fixed coordinates, whereas their entries in a spectral
basis of \(L\) generally depend on \(G\).

In the next identity, \(\mathcal N=x\partial_x-\partial_x^2\) acts on
the scalar profile before evaluation at \(L\). It is not the operator
\(\mathcal N_G\) acting on the resulting Gaussian matrix field.

\begin{lemma}[Gaussian spectral identity]\label{lem:linear-spectral-identity}
Let \(F_*\) be a smooth scalar function whose derivatives have polynomial
growth, and put \(f=F_*''\).  Then
\[
 \mathbb E[(\mathcal N F_*)(L)]
       =(\mathbf S-I)\circ\mathbb E f(L)+\mathfrak R(f).
 \tag{24}\label{eq:24}
\]
\end{lemma}

\begin{proof}
For a smooth scalar function \(h\), its matrix derivative at a real
symmetric matrix satisfies
\[
 \bigl(Dh(L)[M]\bigr)_{ij}
   =\left(\int_0^1 h'((1-t)x_i+tx_j)\,dt\right)M_{ij}.
\]
This is the Dalecki\u{\i}--Kre\u{\i}n divided-difference derivative formula; see
\cite[Theorem~2.11]{Noferini2017}.  To justify it directly here, first use
polynomials and then approximate \(h\) in \(C^1\) on a compact interval
containing the spectrum.  The divided differences converge uniformly, and
the derivative maps converge in Hilbert--Schmidt norm.  This gives the
formula locally, including coincident eigenvalues, without differentiating
a choice of eigenvectors.

Apply this formula to \(h=F_*'\).  In the fixed coordinates,
\[
 \partial_{G_\ell}F_*'(L)=DF_*'(L)[M_\ell]
                         =\bar f\odot M_\ell.
\]
The last expression is interpreted by the preceding spectral convention.
Its Hilbert--Schmidt norm is at most
\(\sup_{|x|\le\|L\|_{\rm op}}|f(x)|\,\|M_\ell\|_{\rm HS}\).
Since \(\|L\|_{\rm op}\le C|G|\), polynomial growth permits Gaussian
integration by parts, for example by inserting smooth cutoffs in \(G\)
and then letting their supports increase.  Thus the local polynomial
approximation is used only to establish the derivative formula; no global
interchange of polynomial approximants and Gaussian expectations is needed.

Finally, \(L\) commutes with \(F_*'(L)\), so
\[
 \begin{aligned}
 \mathbb E[(\mathcal N F_*)(L)]
 &=\mathbb E[L\circ F_*'(L)]-\mathbb E f(L)\\
 &=\sum_\ell M_\ell\circ
           \mathbb E[\bar f\odot M_\ell]-\mathbb E f(L).
 \end{aligned}
\]
Adding and subtracting \(\mathbf S\circ\mathbb E f(L)\) proves
\eqref{eq:24}.
\end{proof}

\subsection{The linear-field estimate}

The normalization~\eqref{eq:8} will cancel the term involving
\(\mathbf S-I\) in Equation~\eqref{eq:24}.  A second estimate controls the
weighted Dirichlet term by positive Gram matrices, and a commutator
estimate accounts for the remaining covariance error.

\begin{proposition}[Linear-field bound]\label{prop:linear-field-bound}
For the normalized field satisfying~\eqref{eq:8}, with
\(t_-I\le T\le t_+I\), and the fixed profiles satisfying
\eqref{eq:2}--\eqref{eq:4} and \(\mathsf C\le-.341\),
\[
 \Delta_{\rm p}-J_\Sigma(u)+.80\|[L,T]\|_2^2
       \le -1.94\tau T^2.
\]
Consequently, with the notation of Equation~\eqref{eq:23},
\[
 \mathcal E\le -.14\tau T^2-\int(W-W_0)\,d\beta\le0.
\]
If \(\mathcal E=0\), then \(T=0\), the measure \(\nu\) is supported on
the diagonal, and \(\beta=0\).
\end{proposition}

\begin{proof}
\smallskip\noindent\emph{The scalar functional.}
Since \(d\) is even, \(\psi_d\) is even, and the differential
identities~\eqref{eq:2} give
\[
 \psi_d''=\mathcal N d=(\mathcal N+2)\mathsf K,
 \qquad
 (\psi_d-v)''=(\mathcal N+2)(\mathsf C-a_R g'').
\]
For any smooth \(F_*\), one has
\((\mathcal N F_*)''=(\mathcal N+2)F_*''\).
Define the even second primitives
\[
 A_K(x)=\int_0^x(x-y)\mathsf K(y)\,dy,\qquad
 A_C(x)=\int_0^x(x-y)\mathsf C(y)\,dy.
\]
They satisfy \(A_K''=\mathsf K\) and \(A_C''=\mathsf C\), and their
derivatives have polynomial growth. The preceding identities show that
\(\psi_d-\mathcal N A_K\) and
\(\psi_d-v-\mathcal N(A_C-a_Rg)\) have zero second derivative.
Since \(v\) and \(g\) are also even, both differences are even and
therefore constant. Thus, for constants \(c_K,c_C\),
\[
 \psi_d=\mathcal N A_K+c_K,\qquad
 \psi_d-v=\mathcal N A_C-a_R\mathcal N g+c_C.
\]
The reference-mean difference
\(\Delta_\Lambda j=\tau_\Lambda(j(L)-(\gamma j)I)\)
annihilates constants. Scalar Gaussian integration by parts gives
\(\gamma(\mathcal N A_K)=\gamma(\mathcal N A_C)
=\gamma(\mathcal N g)=0\), justified by polynomial growth.
We can therefore apply Lemma~\ref{lem:linear-spectral-identity} to
\(A_K\) and \(A_C\). The remaining term is evaluated directly:
\(\mathcal N g=\mathsf K\), so its matrix expectation is
\(-a_R\mathbb E\mathsf K(L)=-a_R\mathbf K\).

The unweighted trace of the first remainder is
\(\tau\mathfrak R(\mathsf K)=\nu[e_K]\).  Applying
Lemma~\ref{lem:linear-spectral-identity} and using trace cyclicity therefore
gives
\[
 \begin{aligned}
 \Delta_{\rm p}
 &=\nu[e_K]+\tau\bigl[(\mathbf S-I)
                   (\mathbf K+T\circ\mathbf C)\bigr]
             -a_R\tau(T\mathbf K)+\tau_T\mathfrak R(\mathsf C)\\
 &=\nu[e_K]+\tau_T\mathfrak R(\mathsf C)
             -\lambda\tau(\mathbf S T^2)
             -\tau\bigl[T^2(-a_R\mathbf C-a_R\lambda T-\lambda I)\bigr]\\
 &\le\nu[e_K]+\tau_T\mathfrak R(\mathsf C)
             -\lambda\tau(\mathbf S T^2)-1.94\tau T^2.
 \end{aligned}
 \tag{25}\label{eq:25}
\]
For the second line use
\(\mathbf K+T\circ\mathbf C=-\lambda T^2\) from~\eqref{eq:8}.
For the last line, the same normalization and
\(\mathbf C\le-.341I\) give
\[
 -\mathbf C-\lambda T
    =\frac{\mathbf H-\mathbf C}{2}
    \ge\frac{m_*+.341}{2}I,
 \qquad
 a_R\frac{m_*+.341}{2}-\lambda=1.94875>1.94.
\]
Pairing this operator inequality with \(T^2\ge0\) is legitimate without
commutation.  In particular, the term involving \(\mathbf S-I\) has been
eliminated without imposing a sign on that matrix.

\smallskip\noindent\emph{The weighted Dirichlet term.}
By linearity in the trace weight, \(J_\Sigma=J_I+J_T\).  Applying Gaussian
integration by parts to the term containing \(L\), and the derivative
formula above to the Dirichlet term, yields
\[
 J_I(u)=\nu\bigl[\overline{|u|^2}-|\bar u|^2\bigr].
\]
For the weighted part, fix a Gaussian input and express \(T\) and every
\(M_\ell\) in the chosen eigenbasis of \(L\); they need not be diagonal.
Expand the products in \(J_T(u)\) and average the summands with outer
indices \(i,k\) interchanged. The scalar kernel multiplying
\((M_\ell)_{ij}(M_\ell)_{jk}T_{ki}\) is
\[
 \frac12\bigl(\overline{|u|^2}_{ij}
                 +\overline{|u|^2}_{kj}\bigr)
          -\bar u_{ij}\cdot\bar u_{kj}.
\]
Set \(u_j=u(x_j)\).  Replacing both copies of the profile by \(u-u_j\)
leaves this kernel unchanged: its constant and linear terms cancel.
The symmetrized kernel of \(\tau_T\mathfrak R(\mathsf C)\) is
\[
 \frac12\bigl(\overline{\mathsf C}_{ij}-\mathsf C_i
              +\overline{\mathsf C}_{kj}-\mathsf C_k\bigr),
 \qquad \mathsf C_i=\mathsf C(x_i).
\]
Swapping \(i,k\) identifies its two contributions after summation;
thus its contraction can also use the single term
\(\overline{\mathsf C}_{ij}-\mathsf C_i\). Define
\[
 b_{ij}^*=\overline{\mathsf C}_{ij}-\mathsf C_i
                         -\overline{|u-u_j|^2}_{ij}.
\]
Subtracting the shifted kernel of \(J_T(u)\) now gives the kernel of
\(\tau_T\mathfrak R(\mathsf C)-J_T(u)\):
\[
 \frac12(b_{ij}^*+b_{kj}^*)
       +(\bar u_{ij}-u_j)\cdot(\bar u_{kj}-u_j).
\]
For fixed \(\ell,j\), this is multiplied by
\(T_{ki}(M_\ell)_{ij}(M_\ell)_{kj}\) and summed over \(i,k\);
we have used symmetry to replace \((M_\ell)_{jk}\) by
\((M_\ell)_{kj}\). Swapping \(i,k\) in half of the first term yields
\[
 \sum_{i,k}T_{ki}(M_\ell)_{ij}(M_\ell)_{kj}
       \frac{b_{ij}^*+b_{kj}^*}{2}
   =\sum_i(TM_\ell)_{ij}(M_\ell)_{ij}b_{ij}^*.
\]
The second term is \(\operatorname{tr}(TQ)\), where
\[
 Q_{ik}=(M_\ell)_{ij}(M_\ell)_{kj}
             (\bar u_{ij}-u_j)\cdot(\bar u_{kj}-u_j).
\]
This is the Gram matrix of the vectors
\((M_\ell)_{ij}(\bar u_{ij}-u_j)\), so \(Q\ge0\) and
\(\operatorname{tr}(TQ)\le t_+\operatorname{tr}Q\).
Sum over \(\ell,j\), take expectation, and divide by \(m\).
The definition of \(\nu\) then gives
\[
 \begin{aligned}
 \tau_T\mathfrak R(\mathsf C)-J_T(u)
 &\le\frac1m\mathbb E\sum_{\ell,i,j}
       (TM_\ell)_{ij}(M_\ell)_{ij}b_{ij}^*
             +t_+\nu[|\bar u_{ij}-u_j|^2]\\
 &\le\lambda\tau(\mathbf S T^2)
          +\frac1{4\lambda}\nu[(b_{ij}^*)^2]
          +t_+\nu[|\bar u_{ij}-u_j|^2].
 \end{aligned}
 \tag{26}\label{eq:26}
\]
The second step is scalar Young's inequality entry by entry, followed by
\(\sum_\ell\|TM_\ell\|_{\rm HS}^2
 =\operatorname{tr}(\mathbf S T^2)\).
The first term in this bound is exactly the one retained with a minus
sign in~\eqref{eq:25}.

\smallskip\noindent\emph{The commutator term.}
For any fixed \(M=M_\ell\), commute \(M\) with the quadratic matrix
equation in~\eqref{eq:8}.  Expanding the Jordan products gives
\[
 \mathbf H\circ[M,T]=[M,\mathbf K]+T\circ[M,\mathbf C].
\]
The matrix \([M,T]\) is skew-symmetric, but the lower norm estimate still
holds.  In an orthonormal basis diagonalizing \(\mathbf H\), the map
\(B\mapsto\mathbf H\circ B\) multiplies entry \(ij\) by
\((h_i+h_j)/2\ge m_*\); hence it increases the Hilbert--Schmidt norm by
at least the factor \(m_*\) on all matrices.  Likewise,
\[
 \|(T-t_cI)\circ B\|_2\le t_r\|B\|_2,
\]
because the spectrum of \(T-t_cI\) lies in \([-t_r,t_r]\).
Splitting the right side at \(t_cI\) and using
\(\|A+B\|_2^2\le1.25\|A\|_2^2+5\|B\|_2^2\) gives
\[
 m_*^2\|[M,T]\|_2^2
 \le1.25\|[M,\mathbf K+t_c\mathbf C]\|_2^2
       +5t_r^2\|[M,\mathbf C]\|_2^2.
\]
All matrices \(M_\ell\) are deterministic in this step.  Thus, in fixed
coordinates,
\[
 [M,\mathbb E f(L)]=\mathbb E[M,f(L)],
 \qquad
 \|\mathbb E[M,f(L)]\|_2^2
      \le\mathbb E\frac1m\|[M,f(L)]\|_{\rm HS}^2.
\]
Only after Jensen's inequality do we pass to an eigenbasis of \(L\).
There the commutator has entries
\((M_\ell)_{ij}(f(x_j)-f(x_i))\).  The resulting endpoint weights are
therefore precisely those defining \(\nu\).  Gaussian orthogonality also
gives \(\|[L,T]\|_2^2=\sum_\ell\|[M_\ell,T]\|_2^2\).
With the definition of \(\Gamma\) in~\eqref{eq:4}, these observations yield
\[
 .80\|[L,T]\|_2^2
       =.80\sum_\ell\|[M_\ell,T]\|_2^2\le\nu[\Gamma].
 \tag{27}\label{eq:27}
\]

\smallskip\noindent\emph{Assembly and strictness.}
Combine~\eqref{eq:25}--\eqref{eq:27}.  The two terms involving
\(\lambda\tau(\mathbf S T^2)\) cancel, giving
\[
 \begin{aligned}
 &\Delta_{\rm p}-J_\Sigma(u)+.80\|[L,T]\|_2^2\\
 &\quad\le\nu\biggl[e_K+\Gamma
       +\frac{(b_{ij}^*)^2}{4\lambda}
       +t_+|\bar u_{ij}-u_j|^2\\
 &\hspace{20mm}
       -\bigl(\overline{|u|^2}_{ij}-|\bar u_{ij}|^2\bigr)\biggr]
       -1.94\tau T^2.
 \end{aligned}
\]
The measure \(\nu\) is symmetric.  Averaging the integrand with the one
obtained by swapping \(i,j\) changes its square term to
\(\bigl((b_{ij}^*)^2+(b_{ji}^*)^2\bigr)/(8\lambda)\), and its
mean-displacement term to
\(t_+(|\bar u_{ij}-u_j|^2+|\bar u_{ij}-u_i|^2)/2\).
In the first square, \(j\) is the endpoint \(\sigma\) of~\eqref{eq:4}
and \(i\) is the opposite endpoint.  The symmetrized integrand is therefore
strictly negative when \(x_i\ne x_j\), by~\eqref{eq:4}, and is zero when
the endpoints coincide.  All endpoint terms are integrable by polynomial
growth, so this proves the linear-field bound in
Proposition~\ref{prop:linear-field-bound}.

Substituting the displayed estimate into Equation~\eqref{eq:23} bounds \(\mathcal E\)
above by the same nonpositive
\(\nu\)-integral, together with
\(-.14\tau T^2-\int(W-W_0)\,d\beta\).
The layer integral is finite by the projection-layer tails and the Gaussian
tails of \(L\), and \(W-W_0>0\) by~\eqref{eq:3} and the averaging
argument preceding~\eqref{eq:22}.  This proves the asserted bound on
\(\mathcal E\), in particular \(\mathcal E\le0\).
If \(\mathcal E=0\), these three nonpositive contributions must all vanish.
It follows that \(T=0\), \(\nu\) assigns zero mass to distinct endpoints,
and \(\beta=0\).  No uniform strict margin is required: a nonnegative
measurable function that is positive on a set can have integral zero only
if that set has measure zero.
\end{proof}

\subsection{Equality and the volume product}

\begin{proof}[Proof of Theorem~\ref{thm:mahler}]
In dimension one, write $K=[a,b]$ with $a<b$. For $a<z<b$,
\[
 |K|\,|(K-z)^\circ|
 =(b-a)\left(\frac1{z-a}+\frac1{b-z}\right)\ge4,
\]
with equality exactly at $z=(a+b)/2$. Every one-dimensional convex
body is a simplex, so this proves the theorem for $n=1$.
For the remaining argument let $n\ge2$ and $m=n+1$.

The Laplace bound~\eqref{eq:1} follows from Equation~\eqref{eq:10} and
Proposition~\ref{prop:linear-field-bound}.  Through the cone-volume identity
this gives
the required lower bound at every translation point
\(z_0\in\operatorname{int}K\).

Suppose first that equality holds in the Laplace bound.  Equation~\eqref{eq:10}
gives \(0\ge-\mathcal E\ge0\), so \(\mathcal E=0\).  Hence \(\nu\)
is supported on the diagonal and \(\beta=0\).  The diagonal support has
the following basis-independent consequence:
\[
 \nu[(x-y)^2]
   =\frac1m\sum_\ell\mathbb E\|[L,M_\ell]\|_{\rm HS}^2
   =\frac1m\sum_{\ell,k}\|[M_k,M_\ell]\|_{\rm HS}^2.
\]
The second identity uses \(L=\sum_kG_kM_k\) and Gaussian orthogonality.
Its left side is zero, so all the deterministic real symmetric matrices
\(M_\ell\) commute.  Choose a single orthonormal basis that diagonalizes
them simultaneously; in this basis \(L\) is diagonal for every \(G\).
Repeated eigenvalues do not affect this conclusion.

In this fixed basis, the vanishing of \(\beta\) says that the
nonnegative integrand
\[
 \sum_i\bigl((P_z)_{ii}-\mathbf1_{\{x_i\le z\}}\bigr)(x_i-z)
\]
vanishes for almost every \((G,z)\).  Away from the spectral endpoints,
each diagonal entry of \(P_z\) is therefore the corresponding zero or
one of \(\Theta_z(L)\).  A positive semidefinite matrix with a zero
diagonal entry has a zero row and column there.  Apply this fact to
\(P_z\) at a zero entry and to \(I-P_z\) at an entry equal to one.
Since \(0\le P_z\le I\), it follows that
\(P_z=\Theta_z(L)\) for almost every \((G,z)\).

For every \(G\) the spectral endpoint set contains only finitely many
values of \(z\), so it has zero product measure.  Fubini's theorem now
provides one fixed layer \(z\) at which
\(D\Pi_C(Z+\xi_z)\) is diagonal almost surely.  The covariance
\(\Sigma\) is positive definite, so \(Z+\xi_z\) has an everywhere
positive Gaussian density.  Thus \(D\Pi_C\) is diagonal Lebesgue-almost
everywhere in the same fixed basis.

The projection \(\Pi_C\) is globally Lipschitz.  Its cross weak
derivatives vanish, and consequently its \(i\)th component depends only
on coordinate \(i\).  For completeness, mollifying gives smooth maps
with the same vanishing cross derivatives; each mollified component is
constant along the other coordinate directions, and local uniform
convergence gives the assertion for \(\Pi_C\).  Write
\[
 \Pi_C(x_1,\ldots,x_m)=(f_1(x_1),\ldots,f_m(x_m)).
\]
Its image is both \(C\) and the Cartesian product of the ranges of the
\(f_i\). Since \(\Pi_C(0)=0\), every \(f_j(0)=0\), and evaluation
at \(t e_i\) gives \(\Pi_C(t e_i)=f_i(t)e_i\in C\).
Conversely, if \(t e_i\in C\), the projection fixes that point, so
\(f_i(t)=t\). The \(i\)th range is therefore exactly the axis section
\(\{t\in\mathbb R:t e_i\in C\}\), a closed convex cone.
Solidity of \(C\) excludes the factor \(\{0\}\), and pointedness
excludes the factor \(\mathbb R\).  Each factor is therefore a half-line,
so the normalized cone is an orthant up to signs.

Undoing the invertible transformation shows that the original cone has
\(m\) linearly independent generating rays.  Every nonzero point of the
original cone has strictly positive height, so each ray meets the
height-one hyperplane. Rescale the generators to these intersection
points. In a nonnegative linear combination of the rescaled generators,
height one means that the coefficients sum to one. The section is
therefore their convex hull. An affine relation among the intersection
points would be a linear relation among the independent generators, so
these \(m=n+1\) points are affinely independent. Thus \(K\) is a simplex.
If equality holds for the infimum defining \(P(K)\), apply this argument
at its interior attaining point.

Conversely, the translated volume product is unchanged by applying an
invertible affine map to both the body and its translation point: the
linear part on the body is paired with its inverse transpose on the polar.
It suffices to take
\[
 K=\operatorname{conv}\{0,e_1,\ldots,e_n\},
 \qquad z_0=\frac{\mathbf1}{m},
\]
where \(\mathbf1\) is the all-ones vector.  The defining polar inequalities
give
\[
 (K-z_0)^\circ
   =\operatorname{conv}\{m\mathbf1,-m e_1,\ldots,-m e_n\}.
\]
Indeed, writing \(s=\sum_{i=1}^n y_i\), the polar inequalities are
\(s\ge-m\) and \(y_i\le1+s/m\). Set
\[
 \alpha_0=\frac{s+m}{m^2},\qquad
 \alpha_i=\alpha_0-\frac{y_i}{m}\quad(1\le i\le n).
\]
These coefficients sum to one and satisfy
\(y=\alpha_0(m\mathbf1)+\sum_{i=1}^n\alpha_i(-m e_i)\).
They are nonnegative exactly when the polar inequalities hold, proving
the displayed convex-hull formula. Its volume is
\[
 |(K-z_0)^\circ|
   =\frac{m^n\det(I+\mathbf1\mathbf1^{\mathsf T})}{n!}
   =\frac{m^m}{n!}.
\]
Since \(|K|=1/n!\), this gives equality in the lower bound.  The scalar
estimates used in the argument are established in
Appendices~\ref{sec:07-scalar} and~\ref{sec:08-segments}.
\end{proof}
